\documentclass[12pt,a4paper]{article}
\usepackage[utf8]{inputenc}
\usepackage[serbian]{babel}
\usepackage{geometry}
\usepackage{longtable}
\usepackage{booktabs}
\usepackage{xcolor}
\usepackage{listings}
\usepackage{adjustbox}
\geometry{left=0.3cm,right=0.3cm,top=1cm,bottom=1cm}

\lstset{
  basicstyle=\ttfamily\footnotesize,
  breaklines=true,
  frame=none,
  numbers=none,
  showstringspaces=false
}

\title{Tabela programskih jezika - Analiza povezivanja varijabli}
\author{}
\date{}

\begin{document}

\maketitle

\begin{longtable}{|p{2cm}|p{5cm}|p{8cm}|}
\hline
\textbf{Jezik} & \textbf{Primjer koda (petlje, nizovi, grananja)} & 
\textbf{Analiza povezivanja varijabli} \\
\hline
\endfirsthead

\hline
\textbf{Jezik} & \textbf{Primjer koda (petlje, nizovi, grananja)} & 
\textbf{Analiza povezivanja varijabli} \\
\hline
\endhead

\hline
\endfoot

\hline
\endlastfoot

\textbf{C} & 
\begin{lstlisting}[language=C]
int arr[10];
int i, sum = 0;
for (i = 0; i < 10; i++) {
    if (arr[i] > 0) {
        sum += arr[i];
    }
}
\end{lstlisting} & 
\textbf{(a) Vidljive varijable, tipovi, vreme života:}

\textit{Lokalne:} \texttt{i, sum, arr} - \texttt{int}, stek-dinamičke. \textit{Globalne:} statičke.

\vspace{0.2cm}
\textbf{(b) Statičko/dinamičko povezivanje:}

Statičko povezivanje tipova (eksplicitna deklaracija). Statički opseg.

\vspace{0.2cm}
\textbf{(c) Karakteristike tipova podataka:}

Primitivni tipovi (\texttt{int, float, char}). Nizovi su statički. Pointeri za dinamičku alokaciju. \\
\hline

\textbf{Java} & 
\begin{lstlisting}[language=Java]
int[] arr = new int[10];
int sum = 0;
for (int i = 0; i < arr.length; i++) {
    if (arr[i] > 0) {
        sum += arr[i];
    }
}
\end{lstlisting} & 
\textbf{(a) Vidljive varijable, tipovi, vreme života:}

\textit{Lokalne:} \texttt{i, sum} - \texttt{int}, stek-dinamičke. \textit{Instance:} \texttt{arr} - referenca na niz, heap-dinamička.

\vspace{0.2cm}
\textbf{(b) Statičko/dinamičko povezivanje:}

Statičko povezivanje tipova. Statički opseg (bez ugniježdenih potprograma). Nema implicitnog povezivanja.

\vspace{0.2cm}
\textbf{(c) Karakteristike tipova podataka:}

Primitivni tipovi i objekti. Nizovi su objekti na heap-u. Automatko upravljanje memorijom (GC). \\
\hline

\textbf{C++} & 
\begin{lstlisting}[language=C++]
int arr[10];
int sum = 0;
for (int i = 0; i < 10; i++) {
    if (arr[i] > 0) {
        sum += arr[i];
    }
}
\end{lstlisting} & 
\textbf{(a) Vidljive varijable, tipovi, vreme života:}

\textit{Lokalne:} \texttt{i, sum, arr} - \texttt{int}, stek-dinamičke. Varijabla \texttt{i} opseg samo u for petlji.

\vspace{0.2cm}
\textbf{(b) Statičko/dinamičko povezivanje:}

Statičko povezivanje tipova. Statički opseg. Blokovi mogu imati lokalne varijable.

\vspace{0.2cm}
\textbf{(c) Karakteristike tipova podataka:}

Primitivni tipovi, klase, template. Nizovi statički ili dinamički (new/delete). RAII za upravljanje resursima. \\
\hline

\textbf{C\#} & 
\begin{lstlisting}[language=CSharp]
int[] arr = new int[10];
int sum = 0;
for (int i = 0; i < arr.Length; i++) {
    if (arr[i] > 0) {
        sum += arr[i];
    }
}
\end{lstlisting} & 
\textbf{(a) Vidljive varijable, tipovi, vreme života:}

\textit{Lokalne:} \texttt{i, sum} - \texttt{int}, stek-dinamičke. \textit{Instance:} \texttt{arr} - referenca, heap-dinamička. Moguće \texttt{var} za zaključivanje tipa.

\vspace{0.2cm}
\textbf{(b) Statičko/dinamičko povezivanje:}

Statičko povezivanje tipova (eksplicitno ili \texttt{var}). Statički opseg. Moguće \texttt{dynamic} za dinamičko povezivanje.

\vspace{0.2cm}
\textbf{(c) Karakteristike tipova podataka:}

Primitivni tipovi i objekti. Nizovi su objekti. Automatko upravljanje memorijom. \\
\hline

\textbf{Python} & 
\begin{lstlisting}[language=Python]
arr = [1, 2, 3, 4, 5]
sum = 0
for i in range(len(arr)):
    if arr[i] > 0:
        sum += arr[i]
\end{lstlisting} & 
\textbf{(a) Vidljive varijable, tipovi, vreme života:}

\textit{Lokalne:} \texttt{i, sum, arr} - implicitno deklarisane, heap-dinamičke. Varijable nemaju tipove, objekti imaju.

\vspace{0.2cm}
\textbf{(b) Statičko/dinamičko povezivanje:}

Dinamičko povezivanje tipova. Statički opseg (dopušta ugniježdene potprograme). Implicitna deklaracija pri dodjeli.

\vspace{0.2cm}
\textbf{(c) Karakteristike tipova podataka:}

Sve je objekat. Dinamički tipovi. Heap-dinamička alokacija. Duck typing. \\
\hline

\textbf{JavaScript} & 
\begin{lstlisting}[language=JavaScript]
var arr = [1, 2, 3, 4, 5];
var sum = 0;
for (var i = 0; i < arr.length; i++) {
    if (arr[i] > 0) {
        sum += arr[i];
    }
}
\end{lstlisting} & 
\textbf{(a) Vidljive varijable, tipovi, vreme života:}

\textit{Lokalne:} \texttt{i, sum, arr} - dinamički tipovi, heap-dinamičke. Varijable mogu mijenjati tip.

\vspace{0.2cm}
\textbf{(b) Statičko/dinamičko povezivanje:}

Dinamičko povezivanje tipova. Statički opseg (dopušta ugniježdene potprograme). Implicitna deklaracija sa \texttt{var}.

\vspace{0.2cm}
\textbf{(c) Karakteristike tipova podataka:}

Sve je objekat ili primitivni tip. Dinamički tipovi. Automatko upravljanje memorijom. \\
\hline

\textbf{PHP} & 
\begin{lstlisting}[language=PHP]
$arr = array(1, 2, 3, 4, 5);
$sum = 0;
for ($i = 0; $i < count($arr); $i++) {
    if ($arr[$i] > 0) {
        $sum += $arr[$i];
    }
}
\end{lstlisting} & 
\textbf{(a) Vidljive varijable, tipovi, vreme života:}

\textit{Lokalne:} \texttt{\$i, \$sum, \$arr} - dinamički tipovi, heap-dinamičke. Globalne varijable nisu direktno vidljive u funkcijama.

\vspace{0.2cm}
\textbf{(b) Statičko/dinamičko povezivanje:}

Dinamičko povezivanje tipova. Statički opseg. Implicitna deklaracija (sve varijable počinju sa \$).

\vspace{0.2cm}
\textbf{(c) Karakteristike tipova podataka:}

Dinamički tipovi. Nizovi su asocijativni. Automatko upravljanje memorijom. \\
\hline

\textbf{Perl} & 
\begin{lstlisting}[language=Perl]
my @arr = (1, 2, 3, 4, 5);
my $sum = 0;
for (my $i = 0; $i < @arr; $i++) {
    if ($arr[$i] > 0) {
        $sum += $arr[$i];
    }
}
\end{lstlisting} & 
\textbf{(a) Vidljive varijable, tipovi, vreme života:}

\textit{Lokalne:} \texttt{\$i, \$sum, @arr} - dinamički tipovi, heap-dinamičke. Implicitna deklaracija po prefiksu (\$, @, \%).

\vspace{0.2cm}
\textbf{(b) Statičko/dinamičko povezivanje:}

Dinamičko povezivanje tipova. Statički opseg. Implicitna deklaracija po prefiksu.

\vspace{0.2cm}
\textbf{(c) Karakteristike tipova podataka:}

Skalarni, nizovi, hash. Dinamički tipovi. Kontekstno određivanje tipa. \\
\hline

\textbf{Ruby} & 
\begin{lstlisting}[language=Ruby]
arr = [1, 2, 3, 4, 5]
sum = 0
arr.each do |i|
    if i > 0
        sum += i
    end
end
\end{lstlisting} & 
\textbf{(a) Vidljive varijable, tipovi, vreme života:}

\textit{Lokalne:} \texttt{i, sum, arr} - reference na objekte, heap-dinamičke. Sve varijable su reference.

\vspace{0.2cm}
\textbf{(b) Statičko/dinamičko povezivanje:}

Dinamičko povezivanje tipova. Statički opseg. Implicitna deklaracija.

\vspace{0.2cm}
\textbf{(c) Karakteristike tipova podataka:}

Sve je objekat. Dinamički tipovi. Duck typing. Automatko upravljanje memorijom. \\
\hline

\textbf{Pascal} & 
\begin{lstlisting}[language=Pascal]
var
    arr: array[1..10] of integer;
    i, sum: integer;
begin
    sum := 0;
    for i := 1 to 10 do
        if arr[i] > 0 then
            sum := sum + arr[i];
end.
\end{lstlisting} & 
\textbf{(a) Vidljive varijable, tipovi, vreme života:}

\textit{Lokalne:} \texttt{i, sum, arr} - \texttt{integer}, stek-dinamičke. Opseg od deklaracije do kraja bloka.

\vspace{0.2cm}
\textbf{(b) Statičko/dinamičko povezivanje:}

Statičko povezivanje tipova (eksplicitna deklaracija). Statički opseg (dopušta ugniježdene potprograme). Nema implicitnog povezivanja.

\vspace{0.2cm}
\textbf{(c) Karakteristike tipova podataka:}

Strogo tipizirani. Nizovi statički. Enumeracijski tipovi. \\
\hline

\textbf{Ada} & 
\begin{lstlisting}[language=Ada]
declare
    type Int_Array is array (1..10) of Integer;
    arr : Int_Array;
    sum : Integer := 0;
begin
    for I in 1..10 loop
        if arr(I) > 0 then
            sum := sum + arr(I);
        end if;
    end loop;
end;
\end{lstlisting} & 
\textbf{(a) Vidljive varijable, tipovi, vreme života:}

\textit{Lokalne:} \texttt{I, sum, arr} - \texttt{Integer}, stek-dinamičke. Varijabla \texttt{I} opseg samo u petlji.

\vspace{0.2cm}
\textbf{(b) Statičko/dinamičko povezivanje:}

Statičko povezivanje tipova. Statički opseg (dopušta ugniježdene potprograme). Nema implicitnog povezivanja.

\vspace{0.2cm}
\textbf{(c) Karakteristike tipova podataka:}

Strogo tipizirani. Nizovi sa eksplicitnim opsegom indeksa. Tipovi podintervala. \\
\hline

\textbf{FORTRAN} & 
\begin{lstlisting}[language=Fortran]
INTEGER :: arr(10)
INTEGER :: i, sum
sum = 0
DO i = 1, 10
    IF (arr(i) > 0) THEN
        sum = sum + arr(i)
    END IF
END DO
\end{lstlisting} & 
\textbf{(a) Vidljive varijable, tipovi, vreme života:}

\textit{Lokalne:} \texttt{i, sum, arr} - \texttt{INTEGER}, statičke ili stek-dinamičke. Implicitna deklaracija (I-N su INTEGER).

\vspace{0.2cm}
\textbf{(b) Statičko/dinamičko povezivanje:}

Statičko povezivanje tipova. Statički opseg. Implicitna deklaracija po prvom slovu.

\vspace{0.2cm}
\textbf{(c) Karakteristike tipova podataka:}

Numerički tipovi (INTEGER, REAL). Nizovi statički. Kompleksni brojevi. \\
\hline

\textbf{BASIC} & 
\begin{lstlisting}[language=Basic]
DIM arr(10)
sum = 0
FOR i = 1 TO 10
    IF arr(i) > 0 THEN
        sum = sum + arr(i)
    END IF
NEXT i
\end{lstlisting} & 
\textbf{(a) Vidljive varijable, tipovi, vreme života:}

\textit{Lokalne:} \texttt{i, sum, arr} - implicitno deklarisane, statičke ili stek-dinamičke. Varijable se deklarišu po prefiksu (\$ za string, \% za integer).

\vspace{0.2cm}
\textbf{(b) Statičko/dinamičko povezivanje:}

Statičko povezivanje tipova. Statički opseg. Implicitna deklaracija po prefiksu ili DIM naredbom.

\vspace{0.2cm}
\textbf{(c) Karakteristike tipova podataka:}

Numerički tipovi i stringovi. Nizovi statički (DIM). Stringovi fiksne ili dinamičke dužine. \\
\hline

\textbf{COBOL} & 
\begin{lstlisting}[language=Cobol]
DATA DIVISION.
WORKING-STORAGE SECTION.
01 arr OCCURS 10 TIMES PIC 9(5).
01 i PIC 9(2).
01 sum PIC 9(8).
PROCEDURE DIVISION.
MOVE 0 TO sum.
PERFORM VARYING i FROM 1 BY 1
    UNTIL i > 10
    IF arr(i) > 0
        ADD arr(i) TO sum
    END-IF
END-PERFORM.
\end{lstlisting} & 
\textbf{(a) Vidljive varijable, tipovi, vreme života:}

\textit{Globalne:} \texttt{arr, i, sum} - deklarisane u DATA DIVISION, statičke. Varijable su globalne za cijeli PROGRAM-ID.

\vspace{0.2cm}
\textbf{(b) Statičko/dinamičko povezivanje:}

Statičko povezivanje tipova (eksplicitna deklaracija u DATA DIVISION). Statički opseg. Nema implicitnog povezivanja.

\vspace{0.2cm}
\textbf{(c) Karakteristike tipova podataka:}

Numerički tipovi (PIC 9), alfanumerički (PIC X). Nizovi (tabele) statički sa OCCURS. Decimalni brojevi s fiksnom tačkom. \\
\hline

\textbf{ALGOL 60} & 
\begin{lstlisting}[language=Algol]
begin
    integer array arr[1:10];
    integer i, sum;
    sum := 0;
    for i := 1 step 1 until 10 do
    begin
        if arr[i] > 0 then
            sum := sum + arr[i]
    end
end
\end{lstlisting} & 
\textbf{(a) Vidljive varijable, tipovi, vreme života:}

\textit{Lokalne:} \texttt{i, sum, arr} - \texttt{integer}, stek-dinamičke. Blokovi stvaraju novi opseg.

\vspace{0.2cm}
\textbf{(b) Statičko/dinamičko povezivanje:}

Statičko povezivanje tipova (eksplicitna deklaracija). Statički opseg (ALGOL 60 je uveo statički opseg). Nema implicitnog povezivanja.

\vspace{0.2cm}
\textbf{(c) Karakteristike tipova podataka:}

Numerički tipovi (integer, real). Nizovi statički sa eksplicitnim opsegom indeksa. Strogo tipizirani. \\
\hline

\textbf{Scheme} & 
\begin{lstlisting}[language=Lisp]
(define (process-array arr)
  (let ((sum 0))
    (do ((i 0 (+ i 1)))
        ((>= i (vector-length arr)))
      (if (> (vector-ref arr i) 0)
          (set! sum (+ sum (vector-ref arr i)))))
    sum))
\end{lstlisting} & 
\textbf{(a) Vidljive varijable, tipovi, vreme života:}

\textit{Lokalne:} \texttt{sum, i} - dinamički tipovi, heap-dinamičke. Varijable u \texttt{let} bloku su lokalne.

\vspace{0.2cm}
\textbf{(b) Statičko/dinamičko povezivanje:}

Dinamičko povezivanje tipova. Statički opseg (dopušta ugniježdene potprograme). Implicitna deklaracija.

\vspace{0.2cm}
\textbf{(c) Karakteristike tipova podataka:}

Sve je objekat. Dinamički tipovi. Vektori (nizovi) i liste. Automatko upravljanje memorijom. \\
\hline

\textbf{Haskell} & 
\begin{lstlisting}[language=Haskell]
processArray arr = 
    sum [x | x <- arr, x > 0]
    
-- ili sa eksplicitnom petljom
processArray arr = 
    let sum = 0
        loop i = if i >= length arr
                 then sum
                 else if arr !! i > 0
                      then loop (i+1) (sum + arr !! i)
                      else loop (i+1) sum
    in loop 0
\end{lstlisting} & 
\textbf{(a) Vidljive varijable, tipovi, vreme života:}

\textit{Lokalne:} \texttt{sum, i} - implicitno tipizirane, heap-dinamičke. Funkcionalni jezik - varijable su nepromjenjive.

\vspace{0.2cm}
\textbf{(b) Statičko/dinamičko povezivanje:}

Statičko povezivanje tipova (zaključivanje tipa). Statički opseg. Implicitno zaključivanje tipa.

\vspace{0.2cm}
\textbf{(c) Karakteristike tipova podataka:}

Strogo tipizirani. Lista je osnovna struktura. Lijena evaluacija. Tipovi klase. \\
\hline

\textbf{ML} & 
\begin{lstlisting}[language=ML]
fun processArray arr =
    let
        val sum = ref 0
        fun loop i =
            if i >= Array.length arr
            then !sum
            else (
                if Array.sub(arr, i) > 0
                then sum := !sum + Array.sub(arr, i)
                else ();
                loop (i + 1)
            )
    in
        loop 0
    end
\end{lstlisting} & 
\textbf{(a) Vidljive varijable, tipovi, vreme života:}

\textit{Lokalne:} \texttt{sum, i} - zaključeni tipovi, heap-dinamičke. \texttt{ref} za promjenjive varijable.

\vspace{0.2cm}
\textbf{(b) Statičko/dinamičko povezivanje:}

Statičko povezivanje tipova (zaključivanje tipa). Statički opseg. Implicitno zaključivanje tipa.

\vspace{0.2cm}
\textbf{(c) Karakteristike tipova podataka:}

Strogo tipizirani. Zaključivanje tipa. Liste i nizovi. Funkcionalni sa bočnim efektima. \\
\hline

\textbf{Lisp} & 
\begin{lstlisting}[language=Lisp]
(defun process-array (arr)
  (let ((sum 0))
    (dotimes (i (length arr))
      (if (> (aref arr i) 0)
          (setf sum (+ sum (aref arr i)))))
    sum))
\end{lstlisting} & 
\textbf{(a) Vidljive varijable, tipovi, vreme života:}

\textit{Lokalne:} \texttt{sum, i} - dinamički tipovi, heap-dinamičke. Varijable u \texttt{let} bloku su lokalne.

\vspace{0.2cm}
\textbf{(b) Statičko/dinamičko povezivanje:}

Dinamičko povezivanje tipova. Dinamički opseg (raniji Lisp) ili statički opseg (Common Lisp). Implicitna deklaracija.

\vspace{0.2cm}
\textbf{(c) Karakteristike tipova podataka:}

Dinamički tipovi. Liste i nizovi. Sve je objekat. Automatko upravljanje memorijom. \\
\hline

\textbf{Modula 2} & 
\begin{lstlisting}[language=Modula2]
MODULE Example;
VAR
    arr: ARRAY [1..10] OF INTEGER;
    i, sum: INTEGER;
BEGIN
    sum := 0;
    FOR i := 1 TO 10 DO
        IF arr[i] > 0 THEN
            sum := sum + arr[i]
        END
    END
END Example.
\end{lstlisting} & 
\textbf{(a) Vidljive varijable, tipovi, vreme života:}

\textit{Lokalne:} \texttt{i, sum, arr} - \texttt{INTEGER}, stek-dinamičke. Opseg od deklaracije do kraja modula.

\vspace{0.2cm}
\textbf{(b) Statičko/dinamičko povezivanje:}

Statičko povezivanje tipova (eksplicitna deklaracija). Statički opseg. Nema implicitnog povezivanja.

\vspace{0.2cm}
\textbf{(c) Karakteristike tipova podataka:}

Strogo tipizirani. Nizovi statički. Moduli za enkapsulaciju. \\
\hline

\textbf{Lua} & 
\begin{lstlisting}[language=Lua]
local arr = {1, 2, 3, 4, 5}
local sum = 0
for i = 1, #arr do
    if arr[i] > 0 then
        sum = sum + arr[i]
    end
end
\end{lstlisting} & 
\textbf{(a) Vidljive varijable, tipovi, vreme života:}

\textit{Lokalne:} \texttt{i, sum, arr} - dinamički tipovi, heap-dinamičke. \texttt{local} za lokalne varijable.

\vspace{0.2cm}
\textbf{(b) Statičko/dinamičko povezivanje:}

Dinamičko povezivanje tipova. Statički opseg. Implicitna deklaracija (bez \texttt{local} su globalne).

\vspace{0.2cm}
\textbf{(c) Karakteristike tipova podataka:}

Dinamički tipovi. Tabele (nizovi i asocijativni nizovi). Sve vrijednosti su objekti prve klase. \\
\hline

\end{longtable}

\end{document}

